\chapter{\MakeUppercase{Conclusão}}\label{sec:conclusao}

Este trabalho caracterizou e ampliou a cobertura de uma representação do
espaço de instâncias do WSP, relacionando propriedades dos casos de teste,
desempenho de seis métodos e trajetórias do IBS. O estudo consolidou e
caracterizou uma base de 55 instâncias por atributos do problema e medidas
topológicas. A geração e a validação de 19 adições para regiões internas
pouco ocupadas do mapa fixo ampliaram a coleção para 74 instâncias.
Compararam-se seis abordagens e analisaram-se quatro famílias de instâncias
nas configurações de ISA. As trajetórias do IBS foram caracterizadas por
redes de soluções visitadas e transições entre elas. Essas contribuições
atendem aos objetivos da Seção~\ref{sec:objetivos}, preservando a distinção
entre o mapa fixo utilizado na geração, os modelos finais de desempenho
e os diferentes protocolos de registro da busca.

A geração orientada demonstrou a viabilidade de produzir instâncias
válidas para regiões internas pouco ocupadas do mapa de referência.
Das 36 buscas, que consumiram 72.000 avaliações, resultaram 19 instâncias
aceitas e distintas. Na união das adições, a maior distância dos pontos
amostrados à coleção diminuiu 20,92\%, e a distância média, 36,27\%.
A ocupação da grade passou de 15\% para 20\%. O preenchimento foi parcial:
52,54\% dos pontos amostrados ainda permaneceram acima do limiar de
subcobertura. Esses resultados respondem ao objetivo de ampliar a cobertura
da representação adotada, mas não demonstram cobertura do conjunto de
instâncias possíveis do WSP nem superioridade sobre uma geração aleatória
com orçamento equivalente.

As 3.936 execuções comparativas mostraram que a ampliação da coleção
revelou regimes de desempenho distintos. O IBS manteve o menor
\textit{gap} médio global, de 3,49\%, e a melhor média, incluindo empates,
em 44 das 74 instâncias. Nas 19 adições, a CUTH apresentou
\textit{gap} médio de 3,23\% e liderou 15 casos. Assim, a liderança global
de um método não substitui a análise dos grupos de instâncias em que
outros são competitivos. Como as adições também alteraram tamanho,
recursos, liberações e tempos de propagação, o estudo não identifica
qual dessas propriedades explica isoladamente a mudança de regime.
Além disso, um \textit{gap} pequeno pode coexistir com perda absoluta
elevada, como ocorreu nas instâncias \textit{Huge} geradas.

Nas duas famílias de instâncias com medidas topológicas adicionais,
quatro dos dez atributos selecionados foram topológicos. Os descritores
escolhidos não coincidiram entre as coleções, e parte deles acompanhou o
tamanho dos grafos. Sua inclusão alterou moderadamente a geometria dos
mapas, que permaneceram semelhantes aos de controle. Nas 74 instâncias,
o diagnóstico auxiliar melhorou
para ILS, LBBD e RS, enquanto CUTH e MIP apresentaram reduções e IBS
permaneceu no mesmo patamar. Na própria amostra, os classificadores do IBS
previram bom desempenho para todos os casos e tiveram acurácia balanceada
de 50\%, apesar da alta acurácia simples. A conclusão sobre topologia é,
portanto, descritiva e condicionada ao método e à configuração utilizada;
a generalização do fluxo completo não foi demonstrada.

As STNs acrescentaram uma perspectiva sobre o percurso da busca que o
objetivo final não fornece. As 160 execuções clássicas e as 190 execuções
nas adições evidenciaram diferenças de recorrência, compartilhamento e
dispersão. Os registros brutos também mostraram retornos a soluções piores
que a incumbente, reforçando a necessidade de distinguir a solução
representativa de cada iteração da melhor solução acumulada. Contudo,
número de nós, compartilhamento e força de entrada não funcionaram como
substitutos diretos da qualidade: nas 19 adições, não se identificou uma
associação monotônica sustentada entre os cinco descritores examinados
e o \textit{gap} médio do IBS. As diferenças de assinatura, parâmetros,
parada e duração efetiva impedem interpretar o contraste entre redes
históricas e adicionais como efeito isolado da expansão.

Em resposta à pergunta da Seção~\ref{sec:pergunta-pesquisa}, foi possível
relacionar perfis estruturais e temporais a grupos de desempenho e
caracterizar comportamentos de busca distintos na coleção ampliada.
A hipótese recebeu sustentação descritiva na identificação do regime
favorável à CUTH nas adições e na caracterização de percursos distintos
de busca. Para a topologia, o resultado se restringe à seleção de
descritores e às alterações moderadas na geometria, com ganhos preditivos
heterogêneos, sem demonstrar informação independente do tamanho.
Não se demonstrou que o aumento de cobertura tenha causado as diferenças
observadas, nem que toda relação detectada dependa da ampliação.

As limitações da Seção~\ref{sec:met-limitacoes} delimitam essas conclusões
à coleção dirigida e aos protocolos avaliados.

Como trabalhos futuros, recomenda-se avaliar o fluxo ISA em instâncias
externas ou por validação aninhada, repetindo o fluxo completo dentro de
cada partição de treinamento e considerando o compartilhamento de
grafos na separação dos dados. Também são prioritárias a sensibilidade
a outros limiares de \textit{gap}, a comparação da geração orientada com
uma geração aleatória sob igual orçamento e a avaliação conjunta da perda
absoluta e do ganho em relação à solução sem recursos. Experimentos que
variem propriedades de forma controlada podem investigar os regimes
favoráveis à CUTH. Para as trajetórias, um protocolo comum de assinatura,
parada e amostragem permitiria comparações mais diretas entre coleções e,
posteriormente, entre métodos.

O contraste central foi a liderança de CUTH em 15 das 19 adições,
incluindo as nove \textit{Huge} e três das seis \textit{Large}, apesar
do menor \textit{gap} médio global do IBS. A média da coleção completa
não evidencia, sozinha, esse regime favorável à CUTH. A ampliação da
base mostrou, assim, que a avaliação dos métodos depende dos perfis
representados e deve considerar seus resultados por grupo, além da
classificação global.
O alcance do estudo é computacional e experimental, sem validar cenários
reais de incêndio nem prescrever decisões operacionais.
